\documentclass[11pt]{article}
\usepackage[a4paper,margin=27mm]{geometry}
\usepackage{amsmath,amssymb,amsthm}
\newtheorem{theorem}{Theorem}
\newtheorem{lemma}[theorem]{Lemma}
\newcommand{\Aut}{\operatorname{Aut}}
\newcommand{\C}{\operatorname{C}}
\title{Enumeration using the product involution}
\author{}\date{}
\begin{document}\maketitle
This note proves the completeness of the algorithm in Section 6 of
Connor and Leemans, \emph{Algorithmic enumeration of regular maps},
Ars Math. Contemp. 10 (2016), 211--222. It does not reproduce the
historical implementation or its running times. The method itself is
attributed to that source.

Let $G$ be a finite group and let $\mathcal I$ be its set of involutions.
Conjugation is $x^g=g^{-1}xg$. Fix $p>2$ and choose one representative
$s$ of every $G$-conjugacy class of elements of order $p$. For each $s$,
choose one representative $\tau$ of each orbit of $\langle s\rangle$
on $\mathcal I$, and put
\[
 r=\tau s^{-1}.
\]
Keep $(r,s)$ if $|r|\ge p$ and $\langle r,s\rangle=G$. Retain one
representative under the simultaneous action of $\Aut(G)$, using the
source's fixed convention for ordered rotations. An explicit further
identification by duality, if desired, must be made separately.

\begin{theorem}
The procedure terminates and gives exactly the automorphism classes
of generating pairs $(r,s)$ with $|s|=p$, $|r|\ge p$ and $|rs|=2$.
\end{theorem}
\begin{proof}
Every retained pair satisfies $rs=\tau$, so it has the required product
order and passes the other two stated tests. Conversely, let $(r_0,s_0)$
satisfy those conditions. Conjugate the pair so that its second entry is
the chosen representative $s$. Its product is then some involution $\tau_0$.
A power $s^i$ sends $\tau_0$ to the chosen representative $\tau$ of its
$\langle s\rangle$-orbit. Conjugating the pair by the same power leaves
$s$ fixed and changes its first entry to $\tau s^{-1}$. Hence a pair in
the required automorphism class is tested and retained unless that class
has already been retained. No class is lost, and the final comparison
removes exactly the duplicates. All lists and orbits are finite.
\end{proof}

Using $\langle s\rangle$ rather than the full centralizer $\C_G(s)$
is sufficient for completeness. It may produce additional candidates;
the final automorphism comparison is still required. Replacing it by
any subgroup of $\C_G(s)$ gives the same completeness statement.
This observation explains the source's cyclic-orbit choice without
assuming that a cyclic orbit is a full stabilizer orbit.

\begin{lemma}
Suppose $G$ has one class of involutions, represented by $\rho$.
The conjugation action on $\mathcal I$ is equivalent to the action on
right cosets $\C_G(\rho)g$, via
\[
 \C_G(\rho)g\longmapsto \rho^g.
\]
The stabilizer of the point labelled $\tau$ is $\C_G(\tau)$.
If $G$ is nonabelian simple, this permutation action is faithful.
\end{lemma}
\begin{proof}
The map is well defined, bijective and equivariant directly from
$\rho^g=\rho^h$ if and only if $gh^{-1}\in\C_G(\rho)$.
An element fixes $\tau$ precisely when it centralizes $\tau$.
The kernel is normal in $G$ and contained in the proper subgroup
$\C_G(\rho)$. Simplicity makes it trivial.
\end{proof}

In the source application, $|G|=460815505920$ and the centralizer of an
involution has order $161280$, so this action has degree $2857239$.
Those group data are explicit imported premises. The algorithm can
carry the involution label with its coset throughout; it need not
recover that label from an abstract stabilizer presentation. In the
source's displayed expression $R=\tau*S^{-1}$, the symbol $S$ must
refer to the current element $s$, not to the sequence of class
representatives previously denoted by the same letter.

For completeness, the source's centralizer-size observation also follows
from the acquired complete subgroup roster. If $g$ has order $m$, then
$\C_G(g)\le N_G(\langle g\rangle)$. A row of class length $n$ in that
roster has normalizer order $|G|/n$. Taking the maximum over \emph{all}
subgroup rows of order $m$ gives the following upper bounds; this does
not require identifying which such rows are cyclic:
\[
\begin{array}{c|rrrrrrrr}
m&3&4&5&6&7&8&10&11\\
\max |N|&6480&161280&720&2160&8232&10752&720&110\\\hline
m&12&14&15&16&19&20&28&31\\
\max |N|&864&588&180&1536&114&80&168&465
\end{array}
\]
Together with the acquired element-order spectrum and unique involution
class, this proves that no nonidentity element has centralizer larger
than $161280$. The involution centralizer attains the bound. The roster
completeness is a borrowed premise, and the calculation uses class
lengths directly; it never multiplies containment weights along chains.

The source also suggests stopping work on a type when its known total
has been reached. This is valid only when that total is exact and the
retained list contains distinct classes satisfying all the tests.
A subset of a finite set with its full cardinality is the whole set.
Consequently the corrected total $9118$ must be used for type
$\{20,31\}$; the source's smaller number $7246$ cannot justify such a
stop. The corrected overall sum is $798644$, replacing $796772$.
These corrections are imported from the accepted numerical treatment,
not proved by the enumeration algorithm itself.

Inversion and the distinction between inner, outer and chiral pairs
are supplied by the accompanying accepted triple correspondence.
The construction above establishes completeness of candidate generation,
not an efficient bound for its runtime. Likewise its polyhedron
interpretation uses the existing intersection-property theorem for
nonabelian simple groups, with the Conder--Oliveros attribution already
recorded in the shared treatment. No new theorem priority is claimed.
\end{document}
